\section{Contoh Program}
\label{sec:contoh-bab7}

Dua contoh berikut memperlihatkan cara kerja format instruksi pada
program nyata: membongkar program dengan \eng{disassembler}, lalu
membandingkan jumlah byte dua bentuk penulisan yang setara.

\subsection{Contoh 1: Membongkar Program dengan Perintah \code{U}}

Program pada Kode~\ref{lst:opcode-demo} sengaja memuat lima
instruksi dengan bentuk operand berbeda, sehingga hasil bongkarnya
dapat dibandingkan dengan kode sumbernya baris demi baris.

\lstinputlisting[
    language=MASM,
    caption={Program \file{opcode-demo.asm} untuk latihan
    \eng{disassembly}.},
    label={lst:opcode-demo}
]{\CodeDir/opcode-demo.asm}

Setelah program dirakit dan dijalankan sekali, program tersebut
dibongkar dengan \file{DEBUG}. Perintah \code{U} dimulai pada alamat
titik masuk program. Hasilnya ditampilkan pada
Kode~\ref{lst:bongkar-opcode-demo}.

\begin{session}[caption={Hasil perintah \code{U} pada
\file{opcode-demo.exe}.}, label={lst:bongkar-opcode-demo}]
-u
0B2A:0000 B82A0B        MOV     AX,0B2A
0B2A:0003 8ED8          MOV     DS,AX
0B2A:0005 B80500        MOV     AX,0005
0B2A:0008 BB0100        MOV     BX,0001
0B2A:000B 03C3          ADD     AX,BX
0B2A:000D 89C1          MOV     CX,AX
0B2A:000F 2D0200        SUB     AX,0002
0B2A:0012 B409          MOV     AH,09
0B2A:0014 8D160000      LEA     DX,[0000]
0B2A:0018 CD21          INT     21
0B2A:001A B44C          MOV     AH,4C
0B2A:001C B000          MOV     AL,00
0B2A:001E CD21          INT     21
\end{session}

Perbandingan kode sumber dan hasil bongkar tersebut dirangkum pada
Tabel~\ref{tab:banding-sumber-bongkar}.

\begin{table}[H]
    \centering
    \caption{Perbandingan kode sumber dan hasil \eng{disassembly}
    pada Kode~\ref{lst:bongkar-opcode-demo}.}
    \label{tab:banding-sumber-bongkar}
    \small
    \begin{tabular}{@{}p{3.6cm}p{2.2cm}cp{4.0cm}@{}}
        \toprule
        \textbf{Kode sumber} & \textbf{Kode mesin} &
        \textbf{Byte} & \textbf{Field yang tampak} \\
        \midrule
        \lstasm|MOV AX, 0005h| & \hex{B8 05 00} & 3 &
        Opcode penambah register, immediate 16 bit \\
        \addlinespace
        \lstasm|MOV BX, 0001h| & \hex{BB 01 00} & 3 &
        Nomor register \reg{BX} = \dec{3}, jadi \hex{B8} $+$ \dec{3} \\
        \addlinespace
        \lstasm|ADD AX, BX| & \hex{03 C3} & 2 &
        ModR/M \bin{11 000 011}: \code{reg} = \reg{AX},
        \code{r/m} = \reg{BX} \\
        \addlinespace
        \lstasm|MOV CX, AX| & \hex{89 C1} & 2 &
        Opcode \hex{89} arah \code{r/m} $\leftarrow$ \code{reg};
        \code{reg} = \reg{AX}, \code{r/m} = \reg{CX} \\
        \addlinespace
        \lstasm|SUB AX, 0002h| & \hex{2D 02 00} & 3 &
        Opcode khusus \reg{AX} dengan immediate 16 bit \\
        \addlinespace
        \lstasm|LEA DX, pesan| & \hex{8D 16 xx xx} & 4 &
        ModR/M \hex{16}: \code{mod} = \bin{00},
        \code{r/m} = \bin{110} (alamat langsung 16 bit) \\
        \bottomrule
    \end{tabular}
\end{table}

\begin{infobox}
Instruksi \lstasm|LEA DX, pesan| tertulis sebagai
\lstasm|LEA DX,[0000]|. Disassembler menampilkan \eng{displacement}
apa adanya, tanpa mengetahui nama simbol \code{pesan}. Nama simbol
hanya hidup di tabel simbol assembler dan hilang seluruhnya pada kode
mesin (Bab~\ref{chap:translasi}).
\end{infobox}

\subsection{Contoh 2: Dua Bentuk Penulisan, Ukuran Berbeda}

Contoh kedua memperlihatkan bahwa dua penulisan yang menghasilkan
keadaan register sama dapat menempati jumlah byte berbeda. Program
pada Kode~\ref{lst:ukuran-instruksi} memuat kedua bentuk tersebut
secara berdampingan dalam satu program \file{.COM}.

\lstinputlisting[
    language=MASM,
    caption={Program \file{ukuran-instruksi.asm} membandingkan dua
    bentuk penulisan.},
    label={lst:ukuran-instruksi}
]{\CodeDir/ukuran-instruksi.asm}

Hasil bongkar kedua bentuk tersebut ditampilkan pada
Kode~\ref{lst:bongkar-ukuran}. Karena program berbentuk \file{.COM},
seluruh instruksi dimulai pada offset \hex{100}.

\begin{session}[caption={Hasil perintah \code{U} pada
\file{ukuran-instruksi.com}, hanya bagian kedua bentuk.},
label={lst:bongkar-ukuran}]
-u 100 10A
13F2:0100 B80100        MOV     AX,0001     ; bentuk A, 3 byte
13F2:0103 BB0100        MOV     BX,0001     ; bentuk A, 3 byte
13F2:0106 BB0100        MOV     BX,0001     ; bentuk B, 3 byte
13F2:0109 89D8          MOV     AX,BX       ; bentuk B, 2 byte
\end{session}

Bentuk A menempati alamat \hex{100} sampai \hex{105}, yaitu enam
byte. Bentuk B menempati alamat \hex{106} sampai \hex{10A}, yaitu
lima byte. Keduanya berakhir dengan \reg{AX} = \dec{1} dan
\reg{BX} = \dec{1}, sehingga hasil akhirnya identik. Selisih satu
byte berasal dari \eng{immediate} 16 bit kedua yang tidak lagi
diperlukan. Rincian selisih tersebut diberikan pada
Tabel~\ref{tab:dua-bentuk}.

\begin{table}[H]
    \centering
    \caption{Perbandingan jumlah byte dua bentuk penulisan yang
    setara.}
    \label{tab:dua-bentuk}
    \small
    \begin{tabular}{@{}p{2.0cm}p{4.6cm}p{3.0cm}c@{}}
        \toprule
        \textbf{Bentuk} & \textbf{Instruksi} &
        \textbf{Kode mesin} & \textbf{Byte} \\
        \midrule
        A & \lstasm|MOV AX, 0001h| & \hex{B8 01 00} & 3 \\
        \addlinespace
        A & \lstasm|MOV BX, 0001h| & \hex{BB 01 00} & 3 \\
        \addlinespace
        \multicolumn{3}{@{}l}{\textbf{Jumlah bentuk A}} & \textbf{6} \\
        \midrule
        B & \lstasm|MOV BX, 0001h| & \hex{BB 01 00} & 3 \\
        \addlinespace
        B & \lstasm|MOV AX, BX| & \hex{89 D8} & 2 \\
        \addlinespace
        \multicolumn{3}{@{}l}{\textbf{Jumlah bentuk B}} & \textbf{5} \\
        \bottomrule
    \end{tabular}
\end{table}

\begin{infobox}
Bentuk B belum tentu lebih baik dalam segala hal. Ia lebih hemat satu
byte, tetapi bergantung pada nilai register lain dan lebih sulit
dibaca. Selisih kecil seperti ini baru menjadi berarti pada rutin
yang dipanggil ribuan kali atau pada memori yang sangat terbatas.
\end{infobox}
